\chapter{2004年全国III卷（文）}
\section{单选题}
\begin{problem}
设集合 \(M = \{(x, y) \mid x^2 + y^2 = 1, x \in \R, y \in \R\}\)，\(N = \{(x, y) \mid x^2 - y = 0, x \in \R, y \in \R\}\)，则集合 \(M \cap N\) 中元素的个数为（\quad）

\choices
    {1}
    {2}
    {3}
    {4}
\end{problem}

\begin{problem}
函数 \( y = \left|\sin \frac{x}{2}\right| \) 的最小正周期是（\quad）

\choices
    {\(\frac{\pi}{2}\)}
    {\(\pi\)}
    {\(2\pi\)}
    {\(4\pi\)}
\end{problem}

\begin{problem}
记函数 \( y = 1 + 3^{-x} \) 的反函数为 \( y = g(x) \)，则 \( g(10) = \)（\quad）

\choices
    {2}
    {\(-2\)}
    {3}
    {\(-1\)}
\end{problem}

\begin{problem}
等比数列 \(\{a_{n}\}\) 中，\(a_{2} = 9, a_{5} = 243\)，则 \(\{a_{n}\}\) 的前4项和为（\quad）

\choices
    {81}
    {120}
    {168}
    {192}
\end{problem}

\begin{problem}
圆 \(x^{2} + y^{2} - 4x = 0\) 在点 \(P(1, \sqrt{3})\) 处的切线方程为（\quad）

\choices
    {\(x + \sqrt{3} y - 2 = 0\)}
    {\(x + \sqrt{3} y - 4 = 0\)}
    {\(x - \sqrt{3} y + 4 = 0\)}
    {\(x - \sqrt{3} y + 2 = 0\)}
\end{problem}

\begin{problem}
\(\left(\sqrt{x} -\frac{1}{x}\right)^6\) 展开式中的常数项为（\quad）

\choices
    {15}
    {\(-15\)}
    {20}
    {\(-20\)}
\end{problem}

\begin{problem}
设复数 \(z\) 的辐角的主值为 \(\frac{2\pi}{3}\)，虚部为 \(\sqrt{3}\)，则 \(z^2 =\)（\quad）

\choices
    {\(-2 - 2\sqrt{3}\i\)}
    {\(-2\sqrt{3} - 2\i\)}
    {\(2 + 2\sqrt{3}\i\)}
    {\(2\sqrt{3} + 2\i\)}
\end{problem}

\begin{problem}
设双曲线的焦点在 \(x\) 轴上，两条渐近线为 \(y = \pm \frac{1}{2} x\)，则该双曲线的离心率 \(e =\)（\quad）

\choices
    {5}
    {\(\sqrt{5}\)}
    {\(\frac{\sqrt{5}}{2}\)}
    {\(\frac{5}{4}\)}
\end{problem}

\begin{problem}
不等式 \(1 < |x + 1| < 3\) 的解集为（\quad）

\choices
    {\((0,2)\)}
    {\((-2,0) \cup (2,4)\)}
    {\((-4,0)\)}
    {\((-4, -2) \cup (0, 2)\)}
\end{problem}

\begin{problem}
正三棱锥的底面边长为 2，侧面均为直角三角形，则此三棱锥的体积为（\quad）

\choices
    {\(\frac{2\sqrt{2}}{3}\)}
    {\(\sqrt{2}\)}
    {\(\frac{\sqrt{2}}{3}\)}
    {\(\frac{4\sqrt{2}}{3}\)}
\end{problem}

\begin{problem}
在 \(\triangle ABC\) 中，\(AB = 3\)，\(BC = \sqrt{13}\)，\(AC = 4\)，则边 \(AC\) 上的高为（\quad）

\choices
    {\(\frac{3\sqrt{2}}{2}\)}
    {\(\frac{3\sqrt{3}}{2}\)}
    {\(\frac{3}{2}\)}
    {\(3\sqrt{3}\)}
\end{problem}

\begin{problem}
将 4 名教师分配到 3 所中学任教，每所中学至少 1 名，则不同的分配方案共有（\quad）

\choices
    {12 种}
    {24 种}
    {36 种}
    {48 种}

\end{problem}

\section{填空题}
\begin{problem}
函数 \( y = \sqrt{\log_{\frac{1}{2}}(x - 1)} \) 的定义域是\fillinblank{}.

\end{problem}

\begin{problem}
用平面 \(\alpha\) 截半径为 \(R\) 的球，如果球心到平面 \(\alpha\) 的距离为 \(\frac{R}{2}\)，那么截得小圆的面积与球的表面积的比值为\fillinblank{}.

\end{problem}

\begin{problem}
函数 \( y = \sin x - \frac{1}{2}\cos x \)（\(x \in \R\)）的最大值为\fillinblank{}.

\end{problem}

\begin{problem}
设 \(P\) 为圆 \(x^{2} + y^{2} = 1\) 上的动点，则点 \(P\) 到直线 \(3x - 4y - 10 = 0\) 的距离的最小值为\fillinblank{}.

\end{problem}

\section{解答题}
\begin{problem}
解方程：\(4^{x} - 2^{x + 2} - 12 = 0\).
\end{problem}

\begin{problem}
已知 \(\alpha\) 为锐角，且 \(\tan \alpha = \frac{1}{2}\)，求 \(\frac{\sin 2\alpha \cos \alpha - \sin \alpha}{\sin 2\alpha \cos 2\alpha}\) 的值.
\end{problem}

\begin{problem}
设数列 \(\{a_{n}\}\) 是公差不为零的等差数列，\(S_n\) 是数列 \(\{a_{n}\}\) 的前 \(n\) 项和，且 \(S_{1}^{2} = 9S_2\)，\(S_4 = 4S_2\)，求数列 \(\{a_{n}\}\) 的通项公式.
\end{problem}

\begin{problem}
某村计划建造一个室内面积为 \(800 \mathrm{~m}^{2}\) 的矩形蔬菜温室. 在温室内，沿左、右两侧与后侧内墙各保留 \(1 \mathrm{~m}\) 宽的通道，沿前侧内墙保留 \(3 \mathrm{~m}\) 宽的空地. 当矩形温室的边长各为多少时？蔬菜的种植面积最大，最大种植面积是多少？
\end{problem}

\begin{problem}
三棱锥 \(P - ABC\) 中，侧面 \(PAC\) 与底面 \(ABC\) 垂直，\(PA = PB = PC = 3\).
\begin{enumerate}
\item 求证：\(AB \perp BC\)；
\item 设 \(AB = BC = 2\sqrt{3}\)，求 \(AC\) 与平面 \(PBC\) 所成角的大小.
\end{enumerate}

\bitmapfigure[max width=0.42\linewidth,max totalheight=4.8cm]{img/2004/national_paper_3_liberal/q21_fig1.png}
\end{problem}

\begin{problem}
设椭圆 \(\frac{x^2}{m+1} + y^2 = 1\) 的两个焦点是 \(F_{1}(-c,0)\) 与 \(F_{2}(c,0)\)（\(c > 0\)），且椭圆上存在一点 \(P\)，使得直线 \(PF_{1}\) 与 \(PF_{2}\) 垂直.
\begin{enumerate}
\item 求实数 \(m\) 的取值范围；
\item 设 \(L\) 是相应于焦点 \(F_{2}\) 的准线，直线 \(PF_{2}\) 与 \(L\) 相交于点 \(Q\). 若 \(\frac{|QF_2|}{|PF_2|} = 2 - \sqrt{3}\)，求直线 \(PF_{2}\) 的方程.
\end{enumerate}
\end{problem}
